\section{综合评价与工程意义}

前文依次研究了单点运输能力、异构多航次调度、连续通信保障以及固定联合任务的独立分区，
并通过敏感性分析和全流程回放对正式方案进行了验证。本章不再重复各问的具体求解过程，
而从统一建模框架、模型与算法特点、核心结果边界以及工程意义四个方面进行综合讨论。

\subsection{四问递进关系与统一建模框架}

四个问题建立在统一的空间、时间、能耗和资源计算规则上，其递进关系可概括为
\[
\text{单点运输与组批}
\rightarrow
\text{多航次运输与资源调度}
\rightarrow
\text{运输--通信协同}
\rightarrow
\text{固定任务独立分区}
\]

问题一在不考虑实体机和共享电池时序的条件下，确定不同服务区的安全运输能力和不可拆货箱组批；
问题二进一步加入多点访问、物资时限、实体无人机和共享电池调度，将能力分析扩展为完整运输时间表；
问题三在相同运输物理规则上增加全过程连续通信约束，并协调运输与中继任务；
问题四则固定问题三的货箱组批、任务顺序、任务时刻及通信关系，仅研究任务分区后各组独立执行的资源代价。

因此，四问并非相互独立，而是逐层增加现实约束。
全文始终采用统一的DEM地形、航段时间、载荷相关能耗、返航安全余量及能源周转规则，
从而保证跨问题结果具有一致的计算口径。

\subsection{模型与算法的主要特点}

首先，本文根据不同问题的结构选择相应求解方法。
问题一的有限组批状态可以完整覆盖，因此采用模式枚举与计数状态动态规划获得限定模型下的精确解；
问题二的路线和资源时序高度耦合，采用``路线列启发式 + 固定路线CP-SAT''的两层框架；
问题三针对离散采样可能遗漏瞬时通信失效的问题，引入连续区间链路证书；
问题四在固定问题三任务后形成有限分区空间，因此对两组和三组方案进行完整枚举。

其次，本文将方案搜索与可执行性验证分离。
正式方案生成后，进一步按照
\[
\text{货箱}
\rightarrow
\text{航次}
\rightarrow
\text{无人机与电池}
\rightarrow
\text{中继资源}
\rightarrow
\text{通信区间}
\]
逐层回放。
其中，问题三共形成4061个连续通信区间，均获得相应的连续可行性证书。
因此，最终结果不仅给出目标函数值，也给出了可核查的执行依据。

最后，本文严格区分不同层级的最优性结论。
问题一由于状态空间被完整覆盖，可给出单点直送模型内的精确最优结论；
问题二和问题三仅对有限候选空间进行搜索与协调，因此正式结果应视为经完整回放验证的高质量可执行方案，
而不扩大为完整组合空间上的全局最优；
问题四则是在固定问题三任务条件下对全部分区进行完整枚举。

\subsection{核心结果汇总与结论边界}

四个问题的核心结果及其适用范围见表~\ref{tab:overall_results}。

\begin{table}[htbp]
    \centering
    \caption{四问核心结果及结论边界}
    \label{tab:overall_results}
    \small
    \begin{tabular}{p{1.0cm}p{4.2cm}p{4.5cm}p{4.7cm}}
        \hline
        问题 & 核心结果 & 主要含义 & 结论边界 \\
        \hline
        一 &
        18架次，59.1313~kWh &
        给出20\%返航安全余量下的单点运输与组批基准 &
        单点直送、固定航迹及本文能耗模型内的词典序精确解 \\

        二 &
        24架次，104.61~min，71.9169~kWh &
        加入多点访问、时限和共享资源后得到完整运输时间表 &
        有限路线搜索后经回放验证的可执行方案 \\

        三 &
        26个运输架次+3个中继架次，
        112.82~min，78.8732~kWh &
        量化连续通信要求带来的额外调度与能源代价 &
        有限运输族和中继窗口内的协调结果，4061个区间全部获证 \\

        四 &
        两组缺3件资源，三组缺7件资源 &
        独立分区削弱共享资源的错峰复用能力 &
        固定问题三任务后的完整分区枚举结果 \\
        \hline
    \end{tabular}
\end{table}

从问题二到问题三，运输架次数由24增加至26，并新增3个中继架次；
联合收尾时间由104.61~min增加至112.82~min，总能耗由71.9169~kWh增加至78.8732~kWh。
这些变化反映了当前候选体系下保证全过程连续通信所需要付出的额外时间、能源和中继资源代价。

问题四进一步表明，统一资源池能够通过跨任务错峰提高利用率。
当任务被划分为相互独立的执行组后，两组和三组方案分别出现3件和7件资源缺口，
说明任务独立性提高会削弱原有资源共享优势。

\subsection{模型局限与改进方向}

本文模型仍存在以下适用边界。

第一，问题二上层路线采用有限启发式搜索，固定路线后的CP-SAT能够严格处理资源与时序约束，
但不能弥补未被路线候选覆盖的组合。后续可采用更大规模邻域搜索、列生成等方法进一步扩充候选空间。

第二，问题三目前采用有限运输方案与有限中继窗口之间的协调，
尚未将通信不可行信息反向反馈至运输主问题。后续可建立运输与通信之间的迭代协调机制，
使通信约束直接参与运输方案生成。

第三，本文采用题给静态DEM和标称设备、传播参数。
现实山区还可能受到强风、降雨、电池衰减、定位误差、植被及多径传播等因素影响。
因此，实际应用时仍需结合实时环境和设备状态重新标定关键参数，
并可进一步建立鲁棒或随机优化模型。

第四，问题四严格固定问题三的联合任务后再进行分区，
因此所得资源缺口表示既定任务被拆分后的独立执行代价。
若允许分区后重新设计运输路线和中继任务，资源需求可能进一步变化。

\subsection{工程启示}

综合四问结果可以看出，山区洪涝条件下的无人机救援并非单纯的最短路径或最少架次问题，
而是地形、载荷、能源、时限、实体资源与通信保障共同作用下的协同调度问题。
实际决策更适合按照
\[
\text{运输能力}
\rightarrow
\text{资源与时序}
\rightarrow
\text{通信保障}
\rightarrow
\text{任务分区}
\]
逐层展开。

同时，通信保障不宜完全作为运输路线确定后的附加环节。
本文从问题二到问题三的结果表明，为满足连续通信要求需要重新协调运输任务并增加中继资源，
因此在山区遮挡明显的环境中，应尽可能将通信可用性提前纳入运输计划。

问题四还表明，任务责任划分虽然有利于独立组织执行，但资源管理不宜过早完全割裂。
若允许无人机、电池或中继设备保持一定程度的跨组共享，可减少重复配置和库存缺口。

最后，第9章敏感性分析表明，当前方案的主要风险集中在少数临界任务：
问题二和问题三最紧运输航次的水平能耗仅允许约0.5441\%的增幅，
问题三最小连续通信裕度仅为0.093780~dB。
因此，本文方案应视为题给标称参数下经过完整回放验证的可执行基准。
实际任务执行前，应结合实时天气、设备状态和实测链路质量重点复核临界航次与最紧通信区间；
若现场条件超过已有安全边界，则应重新进行运输与中继联合调度。